\documentclass[11pt]{article}
% Preámbulo compartido por los artefactos Scrum del proyecto
% Proyecto Agile — Spanish Braille Application (EPN)
\usepackage[utf8]{inputenc}
\usepackage[T1]{fontenc}
\usepackage[spanish]{babel}
\usepackage[a4paper,margin=2.3cm]{geometry}
\usepackage{amsmath}
\usepackage{amssymb}
\usepackage{xcolor}
\usepackage{booktabs}
\usepackage{longtable}
\usepackage{array}
\usepackage{colortbl}
\usepackage{enumitem}
\usepackage{fancyhdr}
\usepackage{titlesec}
\usepackage{tcolorbox}
\usepackage{pgfplots}
\pgfplotsset{compat=1.18}
\usepackage{hyperref}

\definecolor{epnblue}{HTML}{2E4C9C}
\definecolor{epnbluelight}{HTML}{4B6FCB}
\definecolor{epnrow}{HTML}{EDF1FB}
\definecolor{epngray}{HTML}{555555}

\hypersetup{
  colorlinks=true,
  linkcolor=epnblue,
  urlcolor=epnblue,
  citecolor=epnblue,
  pdfcreator={XeLaTeX},
}

\titleformat{\section}
  {\normalfont\Large\bfseries\color{epnblue}}{\thesection}{0.6em}{}
\titleformat{\subsection}
  {\normalfont\large\bfseries\color{epnbluelight}}{\thesubsection}{0.6em}{}

\pagestyle{fancy}
\fancyhf{}
\renewcommand{\headrulewidth}{0.4pt}
\renewcommand{\footrulewidth}{0.2pt}
\fancyfoot[C]{\thepage}

\newcommand{\artefactoTitulo}[2]{%
  \begin{titlepage}
    \centering
    \vspace*{1.2cm}
    {\large Escuela Politécnica Nacional\par}
    {\normalsize Facultad de Ingeniería de Sistemas\par}
    {\normalsize Construcción y Evolución de Software --- Primer Bimestre\par}
    \vspace{1.4cm}
    \begin{tcolorbox}[colback=epnblue,colframe=epnblue,coltext=white,arc=2mm,boxrule=0pt,width=0.92\textwidth,halign=center]
      {\Huge\bfseries #1\par}
    \end{tcolorbox}
    \vspace{0.5cm}
    {\Large #2\par}
    \vspace{1.6cm}
    {\large\bfseries Proyecto Agile --- Spanish Braille Application\par}
    \vspace{0.3cm}
    {\normalsize Traducción bidireccional español \(\leftrightarrow\) Braille Unicode\par}
    \vfill
    \begin{tabular}{ll}
      \textbf{Product Owner:}  & José Castro \\
      \textbf{Scrum Master:}   & Victor Aveiga \\
      \textbf{Equipo de Desarrollo:} & Javier Angulo, Elian Caizapanta, Erick Costa, Emily Aumala \\
    \end{tabular}
    \vspace{0.8cm}
    \par\noindent\rule{0.5\textwidth}{0.4pt}\par
    \vspace{0.3cm}
    {\normalsize Documento preparado por: Erick Costa (Equipo de Desarrollo)\par}
    {\normalsize \today\par}
  \end{titlepage}
  \fancyhead[L]{Proyecto Agile --- Spanish Braille App}
  \fancyhead[R]{#1}
}


\begin{document}
\artefactoTitulo{Incremento}{Producto potencialmente entregable al cierre del Sprint 2}

\section{Definición}
El \textbf{Incremento} es la suma de todos los ítems del Product Backlog completados durante el
Sprint 1 y el Sprint 2, integrados en una aplicación web funcional, ejecutable y potencialmente
entregable al usuario final (personas usuarias de Braille, docentes o instituciones educativas).

\section{Funcionalidad incluida en este Incremento}
\begin{enumerate}[leftmargin=1.4cm]
  \item \textbf{Transcripción Español \(\rightarrow\) Braille} (\texttt{POST /transcribir-Español}):
        convierte abecedario completo, vocales acentuadas, ñ, números con signo numérico, mayúsculas
        (simples y dobles) y signos de puntuación básicos.
  \item \textbf{Transcripción Braille \(\rightarrow\) Español} (\texttt{POST /transcribir-Braille}):
        conversión inversa que reconstruye palabras, números y mayúsculas a partir de una cadena Braille
        Unicode.
  \item \textbf{Braille espejo} (\texttt{POST /espejo}): genera la representación invertida horizontalmente
        (puntos 1\(\leftrightarrow\)4, 2\(\leftrightarrow\)5, 3\(\leftrightarrow\)6) para impresión en
        relieve desde el reverso, con vista dedicada y botón de impresión del navegador.
  \item \textbf{Interfaz web} servida con Thymeleaf: página principal de entrada de texto
        (\texttt{index.html}) y tres vistas de resultado especializadas.
\end{enumerate}

\section{Definition of Done aplicada}
\begin{itemize}[leftmargin=1.4cm]
  \item[\checkmark] Código integrado en la rama principal del repositorio de referencia
        (\texttt{ITERACION-2}).
  \item[\checkmark] Compila y arranca con \texttt{mvn spring-boot:run} / \texttt{mvn -q test}.
  \item[\checkmark] Cobertura de pruebas unitarias para los 3 flujos principales:
        \texttt{BrailleConverterServiceTest} (7 casos: básico, mayúsculas, acentos, ñ, números,
        puntuación, espejo) e \texttt{InverseBrailleServiceTest}.
  \item[\checkmark] Arquitectura refactorizada siguiendo SOLID: un servicio con responsabilidad única
        por regla de negocio (\texttt{AlphabetMapService}, \texttt{AccentMapService},
        \texttt{NumberMapService}, \texttt{UpperCaseHandlerService}, \texttt{PunctuationMapService},
        \texttt{BrailleRendererService}) y un mapa de datos centralizado (\texttt{BrailleCharacterMap}).
  \item[$\sim$] Documentación de diseño actualizada (CRC, casos de uso, diagrama de clases, roadmap).
  \item[$\sim$] Exportación de señalética a PNG/PDF --- el equipo reporta la tarea (Elian Caizapanta)
        como completada, pero no está confirmada en la rama de este repositorio (ver nota en el Sprint
        Backlog); mientras se confirma, se apoya en impresión de navegador como solución provisional.
  \item[\checkmark] \textbf{Este incremento} añade además: contenedor Docker reproducible
        (\texttt{Dockerfile} multi-stage) y documentación de artefactos Scrum en LaTeX/PDF.
\end{itemize}

\section{Cómo ejecutar y validar el Incremento}
\begin{enumerate}[leftmargin=1.4cm]
  \item \textbf{Localmente:} \texttt{./mvnw spring-boot:run} y abrir \url{http://localhost:8080}.
  \item \textbf{Con Docker:} \texttt{docker compose up --build} y abrir \url{http://localhost:8080}
        (ver \texttt{Dockerfile} y \texttt{docker-compose.yml} en la raíz del repositorio).
  \item \textbf{Pruebas automatizadas:} \texttt{./mvnw test} ejecuta la suite JUnit 5 + AssertJ que
        valida los 7 escenarios de transcripción descritos en la Sección 3.
\end{enumerate}

\section{Valor entregado al Product Owner}
Con este incremento, \textbf{José Castro (Product Owner)} cuenta con una aplicación desplegable que
resuelve el objetivo central del producto --- accesibilidad Braille bidireccional en español --- y puede
decidir si el ítem por confirmar (exportación PNG/PDF) se libera en un incremento posterior o si el
producto actual ya aporta valor suficiente para una demostración/entrega parcial (\textit{release}) a
usuarios reales.

\section{Definición de Terminado (Definition of Done)}
Para que una tarea del Sprint Backlog se marque como \textbf{Completada}, el equipo exige que cumpla no
solo ``código escrito'' sino \textbf{``probado e integrado''}: compila sin errores, cuenta con al menos
una prueba unitaria que documenta su comportamiento esperado (ver \texttt{src/test/java}), y su
funcionalidad es accesible desde la interfaz web end-to-end. La suma de todas las tareas que cumplen esta
definición --- no solo las que ``parecen listas'' --- es lo que constituye el \textbf{Incremento}
potencialmente entregable que se demuestra en cada Sprint Review.

\end{document}
